\section{Problem Formulation}
\label{sec:problem}

\paragraph{Setting.}
A robot observes images $o_t$ and proprioceptive states $q_t$.
A frozen visual encoder $\varphi$ maps each image to a grid of patch features.
At a decision time $t$, the context $C_t$ consists of the features of the current and previous images and the recent proprioceptive states; it uses only present and past information.
A candidate action chunk $\bA=(a_t,\dots,a_{t+\He-1})$ is executed for $\He$ steps.
Its \emph{future segment} $\tau$ consists of the frozen features of the images that executing $\bA$ produces, $\varphi(o_{t+1}),\dots,\varphi(o_{t+\He})$, subsampled in time, together with the final proprioceptive state.
A world model must say something about $\tau$ \emph{before} $\bA$ is executed.

\paragraph{Policy-guided planning.}
At every decision, a frozen pretrained policy $\pi$ proposes $K$ chunks $\bA^{(1)},\dots,\bA^{(K)}$ from independent noise seeds, where $\bA^{(1)}$ is the chunk the unmodified policy would execute.
A scorer rates each proposal for the goal.
The highest-scoring proposal is executed, with ties resolved in favor of $\bA^{(1)}$, and the policy replans~\cite{qi2025gpc,cui2026dreamsteer}.
Selecting among the proposals of a competent policy, rather than optimizing actions against the model, keeps candidates on the policy's support.
There, a learned model is less exposed to adversarial exploitation of its errors~\cite{janner2019mbpo,gao2023overoptimization}.
Every compared model receives the same proposal bank.
No scorer can recover an action that is absent from the bank, so we first measure how much a privileged oracle gains by selecting from it (\cref{sec:headroom}).

\paragraph{Goals and labels.}
A goal is specified by a set of goal images $\cG$, such as final frames of successful episodes of the task.
For each candidate, the planner needs to know how far its future segment advances the task toward $\cG$.
The $K$ candidates branch from the same state, so only their \emph{order} matters.
During training, the simulator provides a progress label $\ell_g(\tau)$ for each training goal $g$.
In PushT the label is the fraction of the target region the block covers.
In LIBERO it is a dense progress score built from the goal's BDDL predicates: object-to-target distance for placement predicates and normalized joint progress for articulation predicates.
Labels train the reader only and are never available at test time; held-out goals are never labeled.

\paragraph{Conditional trajectory code.}
A \emph{conditional trajectory code} is a sequence of $M$ discrete tokens $S\in\{1,\dots,V\}^M$ produced by a target encoder $q_\phi(S\mid C,\tau)$, which sees the actual future and is used only during training.
Two heads read the code together with the context, and neither receives $\bA$.
A reader $D_\psi(C,S,g)$ scores progress toward a goal image $g$, and a feature decoder $G_\psi(C,S)$ reconstructs $\tau$.
A world model $p_\theta(S\mid C,\bA)$ predicts the code from the context and the chunk.
Since the heads never see the action, any information about its consequences must pass through $S$.
The code solves a conditional compression problem,
\begin{equation}
\begin{aligned}
  \min_{q_\phi,\,\psi}\ &\E\big[d_{\mathrm{rank}}\big(D_\psi(C,S,g),\ell_g\big)+\alpha\,d_{\mathrm{feat}}\big(G_\psi(C,S),\tau\big)\big]\\
  \text{s.t.}\ &I(\tau;S\mid C)\le R,
\end{aligned}
\label{eq:crd}
\end{equation}
in which the context is available to both the encoder and the heads.
We impose the rate through the code length.
Because $S$ is a deterministic function of $(C,\tau)$, $I(\tau;S\mid C)=H(S\mid C)\le M\log_2 V$, which is 128 bits per segment for $M{=}16$ and $V{=}256$.
This count excludes the codebook, the heads, and the context features, which we report separately.

\paragraph{Predictability.}
At deployment the code is not observed but predicted, $\hat S\sim p_\theta(\cdot\mid C,\bA)$, and what matters is the ranking that $D_\psi(C,\hat S,g)$ induces.
A code that solves~\eqref{eq:crd} may spend bits on aspects of $\tau$ that the action does not determine, such as the outcome of an uncertain contact, and those bits cannot be forecast.
The quantity that matters is the information the action carries about the code, $I(S;\bA\mid C)$.
We estimate it by the difference in cross-entropy between a context-only model $r_\omega(S\mid C)$ and the world model (\cref{sec:wm}).
Optionally, the code is made more predictable during training (\cref{sec:codesign}).

\paragraph{Two limiting cases.}
The formulation contains the two common designs as extremes.
A \emph{per-step latent world model}, such as DINO-WM, sets $S=\tau$.
It applies no compression, predicts the patch features of all $\He$ future steps for every candidate, and needs a readout that is designed or learned on the predicted features.
A \emph{direct scorer} $D(C,\bA,g)$ removes $S$.
Prediction and reading are fused, so every candidate--goal pair needs a full evaluation, and generalization to new goals must be learned jointly with the effect of the action.
\method operates at a finite, explicit rate between them.
We ask whether, at such a rate, predicted codes rank candidates as well as a per-step latent model at lower cost, and whether they transfer to held-out goals better than a direct scorer.
\Cref{tab:designs} summarizes the design properties.

\begin{table}[t]
\centering
\small
\setlength{\tabcolsep}{2.5pt}
\begin{tabular}{@{}llll@{}}
\toprule
 & Goal-independent & Per candidate & Sequential \\
 & prediction & and goal & steps \\
\midrule
Per-step latent WM & $\He\cdot P$ tokens & readout & $\He$ \\
Endpoint latent & $P$ tokens & readout & $1$ \\
Direct scorer & none & full model & $1$ \\
\method (ours) & $M$ tokens & reader & $M$ \\
\bottomrule
\end{tabular}
\caption{\textbf{Design properties} (not results) of models that evaluate a candidate chunk. All four are latent models; $P$ is the number of patch tokens per step ($256$ here), and $M\in\{8,16,32\}$ is the number of code tokens. A direct scorer has no goal-independent part, so its whole cost recurs for every goal. \method's $M$ autoregressive steps are short but not automatically cheaper than $\He$ per-step predictions, so we measure latency rather than assume it.}
\label{tab:designs}
\end{table}
